从一份样本算出一个参数值不难，难的是说清为什么用这条规则而不是另一条，以及凭什么说这一条更好。点估计要同时回答两个问题：怎样把数据变成参数，以及用什么尺子比较两种变法。这两个问题独立——构造原理不自带评价标准，评价标准也不指定构造方法。

先看评价这一头。正态总体抽 \(n=5\) 个点，\(\sigma^2=4\) 已知，比较样本均值 \(\bar X\) 与收缩估计 \(0.6\bar X\)。当真值 \(\mu=1\) 时，\(\bar X\) 无偏、方差 \(0.8\)，MSE 就是 \(0.8\)；收缩估计偏差 \(-0.4\)、方差 \(0.288\)，MSE 是 \(0.448\)。有偏的那个赢了，而且赢得不少。换到 \(\mu=3\)，\(\bar X\) 的 MSE 仍是 \(0.8\)，收缩估计的偏差平方涨到 \(1.44\)，MSE 变成 \(1.728\)，输得也不少。没有哪个估计量在所有参数值上一致最优，所以``哪个更好''这个问题必须先补上``在什么条件下、按什么损失''才有答案。

构造这一头有三条互不相同的原理。矩估计把整份数据压成几个汇总量，让样本矩等于总体矩，再反解参数；它省事，几乎总能算出个数来，代价是丢掉了分布形状里的信息，解还可能落到参数空间外面。最大似然不做压缩，直接问每个候选参数把手上这份数据解释得有多好，取解释力最强的那个；它在正则条件下渐近最优，是从 logistic 回归到语言模型预训练的共同底座，但那句``渐近''是一张要等大样本才兑现的期票。MAP 在似然上乘一个先验再取峰，它和正则化是同一个算式的两种叫法——负对数先验就是罚项，高斯先验给出 \(L_2\)，Laplace 先验给出 \(L_1\)，调罚项系数就是在调先验的强度。

三条路径在有些模型上重合（Bernoulli 的矩估计与最大似然都是 \(\bar X\)），在另一些模型上分道扬镳（均匀分布上矩估计会给出与已观测数据矛盾的参数）。分歧出现的地方，往往正是能看清各自依赖了什么假设的地方。

\subsection{怎么读这一单元}

第一节建立评价语言，后三节按构造原理排开。读的时候不必比较方法的名头，要问三件事：它依赖模型的哪一部分（矩、整条似然、还是外加的先验）；它在优化什么目标（匹配、解释力、还是后验密度）；样本量变化时它的表现怎么走。

四节之间有一条暗线：偏差不是缺陷，而是可以主动买入的东西。第一节用收缩估计说明有偏可以更准，第四节说明这笔交易在贝叶斯语言里就叫先验，中间两节则给出两种不同的``交易入口''。

\subsection{本单元篇目}

以下篇目按概念依赖排列；若从具体问题进入，也可以先打开最接近当前困惑的一篇，再沿篇内链接回补。

\begin{itemize}
\tightlist
\item \hyperref[sec:11-Mathematical-Statistics/02-Point-Estimation/01]{``偏差、方差与均方误差''}；
\item \hyperref[sec:11-Mathematical-Statistics/02-Point-Estimation/02]{``矩估计让样本特征匹配总体特征''}；
\item \hyperref[sec:11-Mathematical-Statistics/02-Point-Estimation/03]{``最大似然选择最能解释数据的参数''}；
\item \hyperref[sec:11-Mathematical-Statistics/02-Point-Estimation/04]{``MAP 把先验与似然合在一个点上''}。
\end{itemize}

四节都停在一个点上。点估计之后还有两个方向：给这个点配一段区间，以及问清楚数据里究竟有多少关于参数的信息——后面的单元分别处理。
